%! Two Categorical Variables — Unit 2, Lesson 2.1 — Homework (Answer Key)
% GENERATED FROM homework/main.tex — the blank with answers filled in; every \answerspace height
% and every \boxguard is identical to it. NO teachernote here — teacher prose lives in the plan.
% GRADUAL RELEASE: this is the lesson's SCORED practice and the LAST component in the packet.
% Paper homework is the DEFAULT for this course; DeltaMath is an occasional override that
% replaces this page and strikes its score cell on the cover. Started in the last eight minutes
% of the period and finished at home; due the first class after two study halls.
% THIRD CONTEXT: the notes teach on phones and sleep, Guided Practice runs on texting drivers,
% and these items are on a gym's membership records — recall is not the skill being built.
% ONE PAGE (compact budget): six items in two parts, then the spiralbox preview of Lesson 2.2.
%   A3 is the denominator trap (of those who STOPPED — divide by the column total, 120).
%   B2 is the crux restated plus the causation hedge; B3 is the AP-style MC with a justification.
% DATA (250 gym members, membership type × still attending after six months):
%              Still attending   Stopped   Total
%   Monthly          60             90      150   -> 40%
%   Annual           70             30      100   -> 70%
%   Total           130            120      250
% \answerspace{H}{} reserves exactly H in the blank; the key passes the answer as the second
% argument, so the two files paginate identically by construction. Keep the macro and every
% \boxguard byte-identical in the blank and the key.
\documentclass[12pt]{article}
\usepackage{apstats-article}
\usepackage{apstats-key}   % pulls in -boxes; do NOT also load -boxes
\newcommand{\answerspace}[2]{\par\nopagebreak\noindent\begin{minipage}[t][#1][t]{\linewidth}%
  \color{keyred}\bfseries #2\end{minipage}\par}

\begin{document}

\pageheader{Unit 2, Lesson 2.1}{Homework --- Answer Key}
% NAMESTRIP: no \namedateperiod here — the name row lives on the cover only.

\vspace{0.06in}

% Authored BYTE-IDENTICALLY in homework/ and homework_key/ — this box is what keeps the
% blank and the key the same number of pages.
\boxguard[8]
\begin{remindbox}
\small
\textbf{This is your graded homework.} It is scored and due the first class after two study
halls. Start it now, ask about anything that stalls you before you leave, then finish it alone.
\end{remindbox}

\vspace{0.04in}

\boxguard[20]
\begin{notesbox}{Part A --- Reading the Table}
\small
\begin{minipage}[c]{0.34\linewidth}\raggedright
A gym checked on its 250 newest members six months after they joined.
\end{minipage}\hfill
\begin{minipage}[c]{0.64\linewidth}
\centering
\renewcommand{\arraystretch}{1.1}
\begin{tabular}{@{} l | c c | c @{}}
 & \textbf{Still attending} & \textbf{Stopped} & \textbf{Total} \\ \hline
\textbf{Monthly plan} & 60 & 90 & 150 \\
\textbf{Annual plan} & 70 & 30 & 100 \\ \hline
\textbf{Total} & 130 & 120 & 250 \\
\end{tabular}
\end{minipage}
\vspace{2pt}
\begin{enumerate}[leftmargin=1.9em, itemsep=2pt]
  \item What proportion of all 250 are annual members still attending? Joint or marginal?
        \answerspace{0.7cm}{$70/250 = 0.28$ --- joint (one cell over the grand total).}
  \item Of the members who \emph{stopped} attending, what proportion had a monthly plan?
        \answerspace{0.7cm}{$90/120 = 0.75$ --- divide by the 120 who stopped.}
\end{enumerate}
\end{notesbox}

\vspace{0.04in}

\boxguard[14]
\begin{notesbox}{Part B --- Is There an Association?}
\small
\begin{enumerate}[leftmargin=1.9em, itemsep=2pt]
  \item Find the proportion still attending for each plan. Show the fractions.
        \answerspace{0.7cm}{Monthly: $60/150 = 0.40$. \quad Annual: $70/100 = 0.70$.}
  \item The manager says: ``70 annual and 60 monthly members still come, so the plan doesn't
        matter.'' Is there an association? Use percents, and say what this study \emph{cannot} show.
        \answerspace{1.3cm}{Yes: the plans differ in size, so compare percents --- 70\% of annual vs.\ 40\% of monthly members still come. Observational data cannot show the plan \emph{causes} it.}
  \item \textbf{AP-style MC.} Which result at another gym shows \textbf{no} association? Circle, then justify.
        \par\vspace{2pt}
        \begin{tabularx}{\linewidth}{@{} X X @{}}
        \textbf{(A)} 60 of 150 monthly, 60 of 100 annual & \ans{(B) 60 of 150 monthly, 40 of 100 annual} \\
        \textbf{(C)} 90 monthly and 70 annual attend & \textbf{(D)} half of all members attend \\
        \textbf{(E)} the two plans have equal sizes & \\
        \end{tabularx}
        \answerspace{0.7cm}{(B): 40\% of each plan still attends --- the conditional distributions match.}
\end{enumerate}
\end{notesbox}

\vspace{0.04in}

% The next-lesson preview closes the packet, so it lives here rather than in AP Practice.
\boxguard[5]
\begin{spiralbox}
\small
\textbf{Next lesson: Scatterplots and Correlation.} Next time both variables are \emph{numbers}
--- hours studied and test score --- and we read their pattern in a scatterplot.
\end{spiralbox}

\end{document}
